\chapter{Entrepreneurship and digital products}
\label{ch:entrepreneurship}

\section{A motivating problem}
A team builds a polished appointment app and nobody adopts it. The failure may be distribution, a weak value proposition, procurement constraints, trust, switching costs, or the fact that the original problem was not costly enough to justify change. Entrepreneurship is disciplined learning under uncertainty, not motivational language.

\learningobjectives{You should be able to formulate value propositions; conduct problem and customer research; segment users and buyers; design minimum viable experiments; reason about metrics, pricing, costs, revenue, distribution, platform effects, accessibility, privacy, and ethical failure; and distinguish product-market evidence from enthusiasm.}

\section{The five-minute explanation}
A product is an offer embedded in a context of users, buyers, alternatives, routines, regulation, and distribution. A value proposition states who has which problem, why existing alternatives are insufficient, and what change the product may create. It is a hypothesis until behaviour or credible evidence supports it.

For CivicQueue, the user may be a citizen, the buyer a municipality, the daily operator a service centre, and the data controller a public organisation. Their incentives and success criteria differ. A product decision must identify whose problem is being tested.

\section{Problem discovery}
Interview for recent behaviour and consequences, not only opinions about a future product. Observe current workarounds. Quantify frequency and cost cautiously. Separate a stated desire from willingness and ability to adopt. Map competitors and non-consumption: a paper form, telephone call, spreadsheet, or doing nothing may be the incumbent alternative.

\section{Minimum viable experiments}
An MVP is not necessarily a small version of the final software. It is the smallest experiment that tests a critical assumption. A concierge workflow, clickable prototype, landing page, manual service, or integration mock can be more informative than a production build.

Write each experiment with:
\begin{enumerate}
\item hypothesis;
\item target participant or buyer;
\item intervention;
\item observable behaviour;
\item threshold or decision rule;
\item duration and cost;
\item ethical and privacy safeguards.
\end{enumerate}
Avoid vanity metrics such as page views when the risk concerns repeated use, payment, completion, or organisational adoption.

\section{Business models and economics}
A business model describes value creation, delivery, and capture. Record fixed costs, variable costs, acquisition costs, support, infrastructure, compliance, and opportunity cost. Pricing can be per user, per transaction, subscription, licence, usage, or public procurement. A price is not proof of value; willingness to pay depends on budget, authority, alternatives, and switching cost.

Platform products face network effects and governance problems. More participants may increase value, but also spam, concentration, lock-in, or unequal visibility. Accessibility and privacy can be competitive and ethical requirements rather than optional features.

\section{Distribution and product-market fit}
Distribution is how a product reaches a context of use. B2B and public-sector distribution may involve procurement frameworks, partnerships, pilots, trust, references, and long sales cycles. Product-market fit is not a single universally valid metric. Look for repeated use, retention, referrals, renewal, low-friction adoption, and evidence that the product solves a meaningful problem for a defined segment.

\section{Ethics and failure}
Growth incentives can reward addictive engagement, surveillance, exclusion, or dark patterns. A product experiment should specify who may be harmed and how participation can stop. A failed experiment is useful when it falsifies an assumption and prevents larger waste. A successful experiment may still be locally valid and not scalable.

\section{Project application}
Include a hypothesis table, interview or observation protocol, experiment log, value proposition, alternative analysis, basic cost model, distribution plan, and ethics section. When presenting the work, explain which assumption is riskiest, what result would cause a pivot, and why the chosen metric represents meaningful behaviour.

\section{Summary and key terms}
Entrepreneurship tests a product system under uncertainty. Discovery begins with problems and alternatives. MVPs test assumptions rather than merely reduce features. Business models include costs, buyers, distribution, regulation, and operations. Metrics should represent behaviour tied to a hypothesis. Ethical design includes what happens when the product succeeds or fails.

\begin{keyterms}problem discovery; value proposition; customer; user; buyer; segment; MVP; hypothesis; experiment; metric; pricing; revenue; cost structure; distribution; product-market fit; platform; network effect; switching cost; ethical design.\end{keyterms}

\section{Exercises and worked solutions}
\exercise{Design a non-software MVP for CivicQueue.}
\solution{A staff member could manually offer a rescheduling form and process requests using a controlled spreadsheet for one service and one week. Measure completion, time, conflicts, citizen understanding, and staff burden. This tests demand and workflow fit before building integration. Personal data would require appropriate safeguards or synthetic records.}

\exercise{Why is a user not always the buyer?}
\solution{A citizen may use a public-service system while a municipality procures it, managers approve it, staff operate it, and another organisation hosts it. Their incentives, risks, and decision rights differ. A product must satisfy the adoption system, not only the end user.}

\exercise{Give one vanity metric and one stronger metric for a rescheduling product.}
\solution{Number of page views is a weak vanity metric. The proportion of eligible users who complete a rescheduling request without staff re-entry, together with error and exclusion rates, is closer to the hypothesised value, though it still requires careful interpretation.}

\section{Further reading}
Treat entrepreneurship as hypothesis testing and connect it to the research methods in Chapter~\ref{ch:research}. A promising experiment is evidence about an assumption, not proof of a business.
